\chapter{Project Roadmap}\label{ch:reg-roadmap}

\section*{Scope of this chapter}
\addcontentsline{toc}{section}{Scope of this chapter}

This chapter is the roadmap for the \textbf{Agentic AI Regulations
corpus and implementation-evidence} project. It is a project-level
artefact, not a programme plan. Other SAP OSS Finance projects (Trial
Balance, Arabic AP, Simula) run on the same platform, workspaces, and
backends, but each owns its own roadmap, business case, and go/no-go
decision. This project is a \emph{supplier} of shared-platform
components (audit service, identity attribution, monitoring,
conformance tooling) that the consumer projects depend on.

\paragraph{Sequence, not dates.}
The milestones and dependency graph below express \emph{order} and
\emph{prerequisite} relationships, not calendar dates. This is
deliberate. Implementation teams must know what depends on what, and
where the schedule can absorb slippage (the Should/Could items); a
date column would harden targets into commitments they were never
meant to be. Dates, when set, live in the project steering packet and
the ticketing system, not in this specification.

\section{MoSCoW classification}\label{sec:reg-moscow}

Every milestone below carries one of four priorities. The mapping is
contractual: a milestone's priority determines what happens when it is
dropped, de-scoped, or slips.

\begin{table}[htbp]
\centering
\caption{MoSCoW prioritisation (shared across the four SAP OSS Finance
projects).}
\label{tab:reg-moscow}
\footnotesize
\begin{tabularx}{\linewidth}{@{}l l X@{}}
\toprule
\textbf{Priority} & \textbf{Meaning} & \textbf{Allowance} \\ \midrule
\textbf{Must} (\textbf{M}) & Non-negotiable for v1.0. Without this milestone, the release does not ship. & None. Missing a Must halts the release; see the kill-switch in \cref{sec:reg-kill-switch}. \\
\textbf{Should} (\textbf{S}) & Important but not vital. Ship without it only with sponsor sign-off; document the deviation in the business-case refresh. & May slip by one consumer wave; two consecutive slips escalate to GFAIC. \\
\textbf{Could} (\textbf{C}) & Desirable enhancement. Included if Must and Should items are on track. & May be dropped to v1.1 at project steering's discretion. \\
\textbf{Won't} (\textbf{W}) & Explicitly out of v1.0 scope. Listed so the boundary is legible. & Not eligible for v1.0 planning at all. \\
\bottomrule
\end{tabularx}
\end{table}

\noindent
The critical path of this project is the chain of Must items in
\cref{tab:reg-milestones}; the Should and Could items are where the
schedule and scope can absorb delivery risk without a release slip.

\section{Dependency graph}\label{sec:reg-dependency-graph}

\begin{figure}[htbp]
\centering
\begin{tikzpicture}[
  node distance=5mm and 10mm,
  must/.style={draw,rounded corners=2pt,align=center,font=\scriptsize,
               fill=regamber!20,draw=regamber!70,
               minimum height=7mm,minimum width=32mm},
  should/.style={draw,rounded corners=2pt,align=center,font=\scriptsize,
               fill=regblue!12,draw=regblue!50,
               minimum height=7mm,minimum width=32mm},
  could/.style={draw,rounded corners=2pt,align=center,font=\scriptsize,
               fill=reggreen!12,draw=reggreen!50,dotted,
               minimum height=7mm,minimum width=32mm},
  dep/.style={-{Latex[length=2mm]},thick},
  soft/.style={-{Latex[length=2mm]},thick,dashed}
]
  \node[must] (m01) {M01 REG-ENG-001 \\ \scriptsize agent index + MGF};
  \node[must,below=of m01] (m02) {M02 REG-AUDIT-001 \\ \scriptsize audit service};
  \node[must,below=of m02] (m03) {M03 REG-IDENT-001 \\ \scriptsize identity attribution};
  \node[must,right=14mm of m02] (m04) {M04 REG-MON-001 \\ \scriptsize monitoring + degradation};
  \node[must,right=14mm of m04] (m05) {M05 REG-INTG-TB-001 \\ \scriptsize TB integration};
  \node[must,below=of m05] (m06) {M06 REG-GA-001 \\ \scriptsize GA sign-off};
  \node[should,right=14mm of m06] (s01) {S01 REG-INTG-AP-001 \\ \scriptsize AP integration};
  \node[should,above=of s01] (s02) {S02 REG-INTG-SIM-001 \\ \scriptsize Simula integration};
  \node[should,above=of s02] (s03) {S03 REG-MGF-V11 \\ \scriptsize MGF v1.1 pillars};
  \node[could,above=of s03] (c01) {C01 REG-V2-EXPLORE \\ \scriptsize multi-regulator};

  \draw[dep] (m01) -- (m02);
  \draw[dep] (m01) -- (m04);
  \draw[dep] (m02) -- (m03);
  \draw[dep] (m03) -- (m05);
  \draw[dep] (m04) -- (m05);
  \draw[dep] (m05) -- (m06);
  \draw[dep] (m06) -- (s01);
  \draw[dep] (s01) -- (s02);
  \draw[dep] (s02) -- (s03);
  \draw[soft] (s03) -- (c01);
\end{tikzpicture}
\caption{Regulations dependency graph. Amber nodes are Must (critical
path); blue are Should (consumer-integration and MGF-pillar
allowance); dotted green are Could (discretionary). Solid arrows are
hard prerequisites; dashed arrows indicate conditional starts (the
target is unlocked when the upstream completes \emph{and} the v2.0
trigger in \cref{sec:reg-v2-trigger} fires).}
\label{fig:reg-dependency-graph}
\end{figure}

\section{Project milestones}\label{sec:reg-milestones}

\begin{table}[htbp]
\centering
\caption{Regulations v1.0 project milestones in dependency order.
Priority: M=Must (critical path), S=Should, C=Could, W=Won't.}
\label{tab:reg-milestones}
\footnotesize
\begin{tabularx}{\linewidth}{@{}l l l l X@{}}
\toprule
\textbf{ID} & \textbf{Prereq} & \textbf{Pri.} & \textbf{Ticket} & \textbf{Milestone} \\ \midrule
M01 & ---       & M & REG-ENG-001        & Agent index and MGF framework specification complete; conformance tooling passes the fixtures in \cref{ch:tooling}. \\
M02 & M01       & M & REG-AUDIT-001      & Audit service (reserve/complete, immutable storage) live in UAT; SEV1/SEV2 clean. \\
M03 & M02       & M & REG-IDENT-001      & Identity attribution wiring (\cref{ch:identity}) live in UAT; REG-MGF-2.1.2-002 compliance upgraded from \textit{missing} to \textit{partial}. \\
M04 & M01       & M & REG-MON-001        & Monitoring spec operational: circuit breaker plus graceful-degradation ladder (\cref{sec:degradation-ladder}). \\
M05 & M03, M04  & M & REG-INTG-TB-001    & First consumer (Trial Balance) integrated end-to-end; reserve/complete coverage $=$ 100\,\% on TB actions. \\
M06 & M05       & M & REG-GA-001         & GA sign-off (audit + identity + monitoring) on the TB integration. \\
M07 & M06       & M & REG-BC-REFRESH     & Business-case refresh and compliance attestation (first cycle after GA; then annual). \\
S01 & M06       & S & REG-INTG-AP-001    & Arabic AP integration: reserve/complete on AP actions; REG-MGF-2.1.2-002 becomes \textit{compliant} once AP live. \\
S02 & S01       & S & REG-INTG-SIM-001   & Simula integration: reserve/complete on Simula corpus-generation actions. \\
S03 & M06       & S & REG-MGF-V11        & MGF v1.1 pillar updates ingested and mapped to agent index. \\
C01 & S03 + \cref{sec:reg-v2-trigger} trigger & C & REG-V2-EXPLORE & v2.0 multi-regulator correlation scoped (beyond IMDA MGF). \\
C02 & M06       & C & REG-TUNE           & Quarterly conformance-tooling retraining from accept/edit telemetry. \\
W01 & ---       & W & REG-V1-REALTIME    & Real-time cross-project audit dashboards are \emph{not} in v1.0 scope. \\
W02 & ---       & W & REG-V1-MULTIREG    & Multi-regulator correlation is \emph{not} in v1.0 scope; single-regulator (IMDA MGF) only. \\
\bottomrule
\end{tabularx}
\end{table}

\noindent
The Must chain (M01 $\to$ M02 $\to$ M03/M04 $\to$ M05 $\to$ M06 $\to$
M07) is the critical path. The Should items (S01--S03) are the
primary \emph{schedule allowance}: additional consumer integrations
and MGF pillars can slip without affecting v1.0 sign-off. The Could
items are the \emph{scope allowance}: they may be cut to v1.1 without
renegotiation.

\section{Launch tiers and exit gates}\label{sec:reg-launch-tiers}

The project walks the four-tier ladder shared across SAP OSS Finance
projects. A tier is exited only when the named exit gate is satisfied;
slippage requires sponsor sign-off.

\begin{table}[htbp]
\centering
\caption{Regulations launch tiers. The exit gate defines the MoSCoW
items that must be closed to exit the tier.}
\label{tab:reg-launch-tiers}
\footnotesize
\begin{tabularx}{\linewidth}{@{}l l X@{}}
\toprule
\textbf{Tier} & \textbf{Scope} & \textbf{Exit gate (references MoSCoW IDs)} \\ \midrule
Alpha & Engineering only, DEV + UAT, conformance fixtures only. & M01, M02, M04 closed; conformance tooling green on fixtures; reproducibility of audit outputs confirmed on two consecutive runs. \\
Beta  & Audit + identity + monitoring live for one consumer project (TB) in UAT. & M03, M05 closed; parallel-run gate in the TB spec passes for one full cycle; no open SEV1/SEV2; audit coverage $=$ 100\,\% of in-scope TB actions. \\
GA (limited) & Audit + identity + monitoring live for TB in production. & M06 closed; Beta gate sustained for two consecutive cycles; kill-switch green; REG-MGF-2.1.2-002 at \textit{partial} (TB-only). \\
Full rollout & All three current consumer projects integrated (TB, AP, Simula). & S01 and S02 closed; REG-MGF-2.1.2-002 at \textit{compliant}; two quarterly reviews sustained; M07 complete; business-case refresh complete. \\
\bottomrule
\end{tabularx}
\end{table}

\section{Kill-switch conditions}\label{sec:reg-kill-switch}

The kill-switch is a pre-committed rulebook: a tripped condition enters
the action column without further debate. GFAIC and sponsor are
notified within 24 hours. The conditions are evaluated continuously
(on every cycle) and do not carry dates.

\begin{table}[htbp]
\centering
\caption{Regulations kill-switch.}
\label{tab:reg-kill-switch}
\footnotesize
\begin{tabularx}{\linewidth}{@{}l l X@{}}
\toprule
\textbf{Condition} & \textbf{Evaluation} & \textbf{Action} \\ \midrule
Audit-coverage gap (missed \texttt{reserve/complete} pair) detected in production & continuous & \textbf{Halt GA onboarding for all consumers}; incident-review before next integration. \\
Circuit breaker tripped with no successful graceful-degradation fallback for three consecutive cycles & per cycle & \textbf{Halt upstream LLM traffic}; fail over to cached/heuristic path; remediate before resuming. \\
Identity-attribution reconciliation fails (principal of record cannot be resolved) for $>$\,0.01\,\% of actions & per cycle & \textbf{Pause affected consumer}; block further reserves until remediated. \\
Any IMDA MGF finding attributable to the Regulations platform (audit, identity, monitoring) & continuous & \textbf{Halt rollout}; remediate before next wave. \\
Any consumer project declares red on a Regulations-supplied component & continuous & \textbf{Pause downstream Should}; re-plan at next steering. \\
\bottomrule
\end{tabularx}
\end{table}

\section{Shared-platform dependencies}\label{sec:reg-shared-deps}

This project depends on two shared-platform components (it supplies
the rest). Each is tracked here only as a dependency, not a commitment
this project can unilaterally adjust. Sequence is expressed in terms
of the Must item that first consumes the dependency, not in calendar
time.

\begin{table}[htbp]
\centering
\caption{Shared-platform dependencies.}
\label{tab:reg-shared-deps}
\footnotesize
\begin{tabularx}{\linewidth}{@{}l l l X@{}}
\toprule
\textbf{Component} & \textbf{Owner} & \textbf{First consumed by} & \textbf{Impact if late} \\ \midrule
Unified schema registry (\path{docs/schema/}) & Shared infra & M01 & Already \texttt{migration.status=complete}; no remaining risk. \\
vLLM hosting (\texttt{vllm-only} policy) & Shared platform & M04 & The graceful-degradation ladder (\cref{sec:degradation-ladder}) is the mitigation; this project is the \emph{owner} of that ladder for sibling consumers. \\
\bottomrule
\end{tabularx}
\end{table}

\noindent
Shared-platform dependencies are surfaced at monthly steering. If any
upstream owner declares \textbf{amber} or \textbf{red} on a component
in \cref{tab:reg-shared-deps}, the affected downstream Must milestone
is re-planned at the following steering.

\section{v2.0 expansion trigger}\label{sec:reg-v2-trigger}

The Could item C01 (REG-V2-EXPLORE) is unlocked only when all of the
following are observed after two consecutive quarters of full rollout
under v1.0:

\begin{itemize}
  \item a second regulator framework (beyond IMDA MGF) is formally
        adopted as in-scope by the AI Governance Working Group;
  \item MGF v1.1 pillars (S03) are fully mapped and at \textit{compliant}
        status;
  \item business-case refresh (M07) explicitly budgets the v2.0
        programme.
\end{itemize}

\noindent
If any of the three conditions fails, C01 remains in the Could column
and is re-evaluated at the next business-case refresh.

\section{Governance and business-case cadence}\label{sec:reg-governance}

\paragraph{Project governance.}
The Regulations project reports monthly to a project steering (AI
Governance Lead, Compliance Officer, Platform Engineer) and quarterly
to the cross-project GFAIC, which it also convenes. Monthly steering
accepts or rejects amber/red signals on Must items; GFAIC reviews
compliance posture and cross-project dependencies (this project's
consumer-integration Should items feed directly into sibling projects'
Must items).

\paragraph{Business-case refresh (M07).}
The business case that authorises this project is refreshed on the
following rhythm:

\begin{itemize}
  \item \textbf{Mandatory refresh}: at the first cycle after GA, and
        annually thereafter (ticket REG-BC-REFRESH, priority M in
        \cref{tab:reg-milestones}). Refresh packet includes consumer
        integration load assumptions, unit economics
        and platform operating cost
        realised vs.\ forecast, and a deviation statement against the
        baseline.
  \item \textbf{Trigger-based refresh}: any of
        \begin{itemize}
          \item consumer-project count in scope changes by $\ge\,1$
                (new consumer adopts, or one is retired);
          \item per-action audit cost changes by $\ge\,30$\,\%
                (storage, indexing, identity-reconciliation throughput);
          \item a new MGF pillar, a new regulator framework, or a
                material change to REG-MGF-2.1.2-002;
          \item any kill-switch in \cref{tab:reg-kill-switch} is
                tripped.
        \end{itemize}
  \item \textbf{Owner}: AI Governance Lead. Approval routed through
        the project steering and filed at
        \path{docs/regulations/business-case/refresh-history.yaml}.
\end{itemize}
